\section{結果：和歌山}
\label{sec:wakayama}

対象は \ang{32.5}--\ang{35.0}N / \ang{134.0}--\ang{137.0}E、
\ang{0.375} 格子で 72 点。紀伊水道から潮岬の南の太平洋までを含む。

北海道と違い、\texttt{sites.json} に和歌山の打ち上げ地点は登録されていない。
そのため本章は\textbf{領域全体の統計}で述べ、
代表地点の表は作っていない。
地点を決めずに 1 点だけ数値を出すと、根拠のない精度を持ち込むことになるためである。
図 \ref{fig:m8-wakayama} だけは、比較のために領域の中央格子点を使っている。

\subsection{年間を通した姿}

\begin{figure}[H]
\centering
\includegraphics[width=0.95\linewidth]{wakayama/F0_sample_counts.pdf}
\caption[標本数]{各ビンに入った解析値の数。前半は必ず 15 日ぶんで 600 個、後半は月の長さで変わり 1 月は 640 個、2 月は閏年が 3 回あるので 532 個になる。この図が合わなければ取得か集計が壊れているので、他のどの図よりも先に確認する。}
\label{fig:f0-wakayama}
\end{figure}

\begin{figure}[H]
\centering
\includegraphics[width=1.0\linewidth]{wakayama/F1_alt_half_speed.pdf}
\caption[時期×高度の風速と風向]{上：平均風速。下：平均風向（吹いてくる向き）。風向には円環用のカラーマップを使っている。順次カラーマップを使うと \ang{0} と \ang{360} の境目に実在しない断層が現れるためである。灰色は定常度が 0.3 未満で塗りを外した領域で、そのビンではそもそも風向が定まっていない。}
\label{fig:f1-wakayama}
\end{figure}

\begin{figure}[H]
\centering
\includegraphics[width=1.0\linewidth]{wakayama/F3_profiles_by_month.pdf}
\caption[月別の高度プロファイル]{月ごとの風速プロファイル。線は各ビンの中央値、灰色の帯はその月の 2 ビンを通した p10--p90 の広がり。帯の幅が平均からのずれの大きさを表しており、線だけを見て計画を立ててはいけない。}
\label{fig:f3-wakayama}
\end{figure}


和歌山の特徴は\textbf{冬のジェットの強さ}である。
\SI{200}{\hecto\pascal}（約 \SI{12.2}{\kilo\metre}）の平均風速は冬に
\SI{75.5}{\metre\per\second}、定常度は 0.98 に達する。
北海道の同じ高度・同じ季節（\SI{40.2}{\metre\per\second}）のほぼ 2 倍である。

これは和歌山が亜熱帯ジェットの真下にあたるためで、
冬に南下した東アジアジェットの核がちょうどこの緯度帯を通る。
定常度 0.98 という値は、
\textbf{冬の 3 か月間、この高度の風はほとんど常に西からの強風だった}ことを意味する。

夏はこれが大きく変わり、\SI{200}{\hecto\pascal} で \SI{22.7}{\metre\per\second}、
定常度 0.57 まで落ちる。ジェットの核が北へ抜けるためである。

\subsection{成層圏}

\begin{figure}[H]
\centering
\includegraphics[width=1.0\linewidth]{wakayama/F2_alt_half_constancy.pdf}
\caption[定常度]{定常度（ベクトル平均風速 ÷ スカラー平均風速）。1 に近いほど風向が安定し、0 に近いほどビンのなかで風向が入れ替わっている。淡い帯が成層圏の風の反転している高度で、季節とともに上下する。}
\label{fig:f2-wakayama}
\end{figure}

\begin{figure}[H]
\centering
\includegraphics[width=1.0\linewidth]{wakayama/F7-1_wind_roses.pdf}
\caption[風配図（前半）]{代表格子点の風配図、\textbf{月の前半}（1--15 日）。
平均と楕円では見せられない「分布そのもの」を示す。
ローブが 2 つに割れている場面が、平均風がほぼゼロで分散だけが大きくなる時期にあたる。}
\label{fig:f7-wakayama-h1}
\end{figure}

\begin{figure}[H]
\centering
\includegraphics[width=1.0\linewidth]{wakayama/F7-2_wind_roses.pdf}
\caption[風配図（後半）]{同じ格子点の\textbf{月の後半}（16 日--月末）。
前半（図\ref{fig:f7-wakayama-h1}）と並べると、ほとんどの月で\textbf{2 枚はよく似ている}。
唯一 10 月だけ \SI{30}{\hecto\pascal} の最頻セクターが前半 \ang{90}、
後半 \ang{248} と入れ替わるが、どちらも頻度は 2 割前後しかない。
つまり向きが変わったというより、\textbf{その時期はどちらの半月も
風向が定まっていない}のであって、最頻セクターを読む意味が薄い。}
\label{fig:f7-wakayama-h2}
\end{figure}

\SI{10}{\hecto\pascal}（約 \SI{31}{\kilo\metre}）は、
夏に \SI{18.2}{\metre\per\second}・定常度 0.99 の\textbf{東風}である。
これは熱帯・亜熱帯の夏の成層圏東風で、
北海道の同じ高度（\SI{13.6}{\metre\per\second}）より強い。

一方、冬と春の定常度はそれぞれ 0.40、0.37 しかない。
定常度が最も低いのは 3 月後半（0.05）、11 月前半（0.11）、
1 月後半（0.18）、3 月前半（0.21）、2 月前半（0.22）である。
北海道と同じく、春と秋の東西逆転期に風向が定まらなくなる。

\subsection{領域の中で答えが変わるか}

\begin{figure}[H]
\centering
\includegraphics[width=1.0\linewidth]{wakayama/F6_spatial_spread.pdf}
\caption[領域内のばらつき]{領域内での平均風速の幅（最大 $-$ 最小）。これが小さい高度・時期では、領域のどこで上げても風はほぼ同じで、代表 1 地点のプロファイルで語ってよい。大きいところではそうではない。}
\label{fig:f6-wakayama}
\end{figure}


領域内での平均風速の幅は最大でも \SI{12.0}{\metre\per\second} である。
北海道の \SI{31.7}{\metre\per\second} と比べるとはるかに小さい。

つまり\textbf{和歌山の領域は、代表 1 地点でおおむね語ってよい}。
領域が狭いこと（北海道の 8 分の 1 の格子点数）に加えて、
緯度の幅が小さくジェットの勾配をまたがないためである。

\subsection{地図で見る}

\begin{figure}[H]
\centering
\includegraphics[width=1.0\linewidth]{wakayama/M1-200_month_grid_200hPa.pdf}
\caption[200 hPa の月別分布図]{\SI{200}{\hecto\pascal}（約 \SI{12}{\kilo\metre}）の平均風を月ごとに並べたもの。塗りが風速、矢印が風向。ジェットが季節とともに動く様子が地図の上で見える。}
\label{fig:m1200-wakayama}
\end{figure}

\begin{figure}[H]
\centering
\includegraphics[width=1.0\linewidth]{wakayama/M1-030_month_grid_30hPa.pdf}
\caption[30 hPa の月別分布図]{同じく \SI{30}{\hecto\pascal}（約 \SI{24}{\kilo\metre}）。対流圏とは逆の季節変化をしており、夏に東風が強まる。}
\label{fig:m1030-wakayama}
\end{figure}

\begin{figure}[H]
\centering
\includegraphics[width=1.0\linewidth]{wakayama/M3-200_direction_map_200hPa.pdf}
\caption[風向の分布図]{\SI{200}{\hecto\pascal} の平均風向。円環用のカラーマップを使い、定常度 0.3 未満は灰色にしている。灰色の領域では平均風向に意味がない。}
\label{fig:m3200-wakayama}
\end{figure}

\begin{figure}[H]
\centering
\includegraphics[width=1.0\linewidth]{wakayama/M2_constancy_map.pdf}
\caption[定常度の分布図]{高度と季節ごとの定常度。濃いところでは風向が安定していて計画が立てやすく、淡いところではビンのなかで風向が入れ替わる。}
\label{fig:m2-wakayama}
\end{figure}

\begin{figure}[H]
\centering
\includegraphics[width=0.92\linewidth]{wakayama/M4-01_wind_ellipse_map.pdf}
\caption[風のばらつきを地図に]{1 月前半、\SI{200}{\hecto\pascal} の平均風ベクトルと、その先端に描いた $1\sigma$ の散らばり楕円。丸い楕円はどの向きにもばらつくこと、細長い楕円は主に強さがばらつくことを意味する。}
\label{fig:m401-wakayama}
\end{figure}


\subsection{落下位置}

\begin{figure}[H]
\centering
\includegraphics[width=1.0\linewidth]{wakayama/M5_drift_month_grid.pdf}
\caption[月別の平均落下位置]{各格子点から上げた場合の平均落下位置。\textbf{矢印は実寸}で、矢先がそのまま予測落下点である。色が平均の飛行距離。全パネルが同じ枠と同じ色スケールなので、月どうしを直接比べられる。}
\label{fig:m5-wakayama}
\end{figure}

\begin{figure}[H]
\centering
\includegraphics[width=0.92\linewidth]{wakayama/M7-01_drift_map.pdf}
\caption[1 月の落下位置]{1 月前半の平均落下位置を拡大したもの。枠は集計対象の領域で、矢先はその外へ出ていることが多い。}
\label{fig:m701-wakayama}
\end{figure}

\begin{figure}[H]
\centering
\includegraphics[width=0.92\linewidth]{wakayama/M7-07_drift_map.pdf}
\caption[7 月の落下位置]{同じく 7 月前半。1 月と比べると距離も向きも大きく違う。}
\label{fig:m707-wakayama}
\end{figure}

\begin{figure}[H]
\centering
\includegraphics[width=1.0\linewidth]{wakayama/M6_drift_percentile_map.pdf}
\caption[荒れた日の飛行距離]{飛行距離の 90 パーセンタイル。平均 \SI{150}{\kilo\metre} で p90 が \SI{400}{\kilo\metre} という状況と、平均 \SI{150}{\kilo\metre} だけを見て計画することは、運用上まったく違う。}
\label{fig:m6-wakayama}
\end{figure}

\begin{figure}[H]
\centering
\includegraphics[width=0.92\linewidth]{wakayama/M8_landing_ellipse_map.pdf}
\caption[代表地点からの落下分布]{1 地点から上げた場合の落下位置を季節別に重ねたもの。楕円は 50 / 90 / 95 \% の範囲で、各解析時刻ごとに求めた着地点の共分散から描いている。散らばりが 2 変量正規に従うという仮定は残る。灰色が陸、白が海である。}
\label{fig:m8-wakayama}
\end{figure}


飛行時間は約 \num{137} 分、領域平均の飛行距離は年間平均 \SI{158}{\kilo\metre}。
北海道（\SI{143}{\kilo\metre}）よりわずかに長い。

\begin{table}[H]
\centering
\caption{和歌山：飛行距離が最短・最長のビン（領域平均）。
方位は機体が\textbf{進む}向き。}
\begin{tabular}{lrrr}
\toprule
時期 & 平均距離 [\si{\kilo\metre}] & 方位 [度] & p90 [\si{\kilo\metre}] \\
\midrule
8 月 後半 & 49 & 285（西北西） & 83 \\
9 月 前半 & 53 & 44（北東） & 87 \\
8 月 前半 & 54 & 259（西南西） & 87 \\
\midrule
2 月 前半 & 276 & 87（東） & 343 \\
1 月 後半 & 270 & 88（東） & 341 \\
12 月 後半 & 269 & 85（東） & 333 \\
\bottomrule
\end{tabular}
\end{table}

北海道と最も大きく違うのは\textbf{夏の飛行方位}である。

\begin{itemize}
  \item 冬（12--2 月）は東へ \SIrange{251}{276}{\kilo\metre}。
        図 \ref{fig:m8-wakayama} のとおり、和歌山中央から上げると
        平均落下点は御前崎の南の海上になる。
  \item 7 月後半から 8 月は\textbf{西寄り}で、方位は \ang{247}--\ang{285}。
        一方で 7 月前半は東（\ang{96}）、9 月前半は北東（\ang{44}）と、
        半月ごとに向きが入れ替わる。
        対流圏上部の弱い西風と成層圏の東風がほぼ釣り合い、
        どちらが少し勝つかで向きが決まるためである。
        北海道の夏（東南東）とは向きが違い、しかも距離が短い。
  \item 8 月後半が年間で最も扱いやすく、平均 \SI{49}{\kilo\metre}、p90 でも \SI{83}{\kilo\metre} に収まる。
\end{itemize}

ただし夏は、東向きの日と西向きの日が打ち消し合うので、
\textbf{平均落下点は打ち上げ地点のすぐ近くになる}。
8 月後半の領域平均は 1 回あたり \SI{49}{\kilo\metre} 飛ぶのに、
平均落下点は \SI{18}{\kilo\metre} しか動かない。
夏の計画は平均の矢印ではなく、散らばり（図 \ref{fig:m8-wakayama} の楕円）で立てる必要がある。
海に落ちやすい北海道とは、回収計画の立て方が大きく変わる。

\subsection{日変化と鉛直シアー}

\begin{figure}[H]
\centering
\includegraphics[width=0.75\linewidth]{wakayama/F8_diurnal.pdf}
\caption[日変化]{夜（03 JST）と昼（15 JST）の平均風速の差を、季節ごとに高度に対して描いたもの
（\SI{300}{\hecto\pascal} 以下の層のみ、領域平均）。横軸の目盛りに注意してほしい。
差は 3 領域・4 季節のどれでも \SI{1}{\metre\per\second} 未満で、
日々の変動（数十 \si{\metre\per\second}）に比べて小さい。
これが\textbf{時刻を区別せず平均してよい根拠}である。
ビンを時刻で 4 分割すると標本数も 4 分の 1 になるので、
必要のないところで分割しない理由がここにある。}
\label{fig:f8-wakayama}
\end{figure}

\begin{figure}[H]
\centering
\includegraphics[width=1.0\linewidth]{wakayama/F9_shear.pdf}
\caption[鉛直シアー]{隣り合う気圧面のあいだの風速差を高度差で割ったもの。
機体そのものはシアーで壊れないが、大きいところでは
高度の推定が少しずれるだけで水平の予測が大きく動く。
最大はジェットの上下と地表付近に出る。}
\label{fig:f9-wakayama}
\end{figure}


夜と昼の平均風速の差は最大 \SI{0.9}{\metre\per\second} で、
北海道と違って高度 \SI{1}{\kilo\metre} 付近（夜のほうが強い）に山がある。
その上では \SI{0.6}{\metre\per\second} 以下に収まる。

鉛直シアーは北海道より大きく、最大 \SI{9.3}{(\metre\per\second)\per\kilo\metre}
（12 月後半、高度 \SI{16.3}{\kilo\metre}）。
地表付近ではなく\textbf{ジェットの直上}が最大になるのが和歌山の特徴で、
冬の亜熱帯ジェットが強く、その上で急に弱まるためである。
この高度帯は破裂高度のすぐ下にあたるので、
高度の推定が数百メートルずれると水平の予測が大きく動く。

\subsection{既存シミュレータの ISA 仮定}

\begin{figure}[H]
\centering
\includegraphics[width=1.0\linewidth]{wakayama/F10_isa_check.pdf}
\caption[ISA との比較]{気圧から高度への換算。左は JRA-3Q の実際の気圧面高度と、既存シミュレータが使う国際標準大気 ISA。右はその差。ISA は季節も緯度も天気も知らないので、季節によって系統的にずれる。}
\label{fig:f10-wakayama}
\end{figure}


ISA からのずれは最大 \SI{787}{\metre}（夏）、冬は \SI{242}{\metre}。
北海道（夏 \SI{1036}{\metre}）より小さいが、やはり無視できる量ではない。
